\begin{lemma}
For every extended nonnegative real $0<d\leq1$ and integer $D\geq0$ there
are $0<\varepsilon<1$, $1\leq C<\infty$ and an integer $N\geq1$, all
chosen before the number of clusters, with this property. For every $r$
and $n\geq N$, let $B$ be a finite blow-up instance with candidacy system
$A$, reduced graph $R$ on $[r]$ of maximum degree at most $D$, and
$(\varepsilon,d)$-super-regular host and candidacy pairs. Suppose each
source cluster has size $n$ or $n+1$ and has the same size as its host
cluster. Suppose every source pair corresponding to a reduced edge is a
saturating matching. Then there is a law on compatible full Copies from
the source to the host whose underlying maps are $(C/n)$-vertex-spread
on the entire host vertex set.
\end{lemma}
\begin{proof}
Set $Q=D^2+1$. Apply \noderef{EqualSizeParameterTable} at $(d,D)$ and
choose its sequences $p,\rho,e,\omega,\theta,\zeta,\widehat\zeta$,
certificate $\mathcal K$, numbers $a,M,C$, and cutoff $N_0$. These choices
precede $r$. The certificate has precisely the meaning in
\noderef{InductionApplicationCertificate}; exposures and rounds remain
independent indices. For $s<Q$, the table gives
$0<e_s\leq e_{s+1}\leq1$. In particular $e_0$ is positive and finite.
Take $\varepsilon=e_0/2$, so $0<\varepsilon<1$ and
$\varepsilon\leq e_0$. Also $1\leq M$ and $2M\leq C$, whence $1\leq C$.
Finiteness of $C$ must not be assumed from notation: take
$k=\max\{N_0,1\}$. The uniform-size conclusion gives $C/k\leq1$.
Since $k$ is positive and finite, cancellation yields $C\leq k<\infty$.
It follows from $2M\leq C$ that $M$ is finite too.

The finite nonempty family $(e_s)_{s<Q}$ has a positive minimum $e_*$.
Choose $N_e\geq1$ with $1/N_e\leq e_*$ and set
$N=\max\{N_0,N_e,1\}$. Thus $n\geq N$ implies
$1/n\leq e_s$ for every $s<Q$. The table's other conclusions remain
valid at every such $n$. In particular $C/n\leq1$ and
$D(1-a)^n\leq1/2$ with $0<a\leq1$.

Fix $r,n,B,A$ as in the statement. Increasing the input error from
$\varepsilon$ to $e_0$ preserves super-regularity. Indeed the tests in
\noderef{RegularPair} become more restrictive while both density bounds
become weaker, and each degree interval in \noderef{SuperRegularPair}
widens. Applying this separately in \noderef{InstanceSuperRegular} and
\noderef{CandidacySuperRegular} gives the input at error $e_0$ needed below.
We check every numerical premise of \noderef{VariableSizeInductionIteration}.
The table gives $p_0=d$, $p_{x+1}=\rho_x(d/2)$ for $x<D$, and the
probability and failure bounds just stated. For $s<Q$ it gives
$e_s\leq e_{s+1}\leq1$, $\omega_{s+1}<e_{s+1}$, and
\[
 2\sqrt{\zeta_s}\leq\theta_s,\quad
 \zeta_s+D\sqrt{\zeta_s}\leq\theta_s,\quad
 \zeta_s+2D\zeta_s\leq\theta_s.
\]
For every $x\leq D$ it gives $0<\theta_s$, $4\theta_s<\rho_x$ and
$2\zeta_s<\rho_x,1-\rho_x$, and it also gives
$2\zeta_s<d/2,1-d/2$.
For the additive density gap, fix $x<D$. Both $p_x,p_{x+1}$ lie in
$(0,1]$. The table gives $2e_s<p_x-p_{x+1}$, forcing
$p_{x+1}<p_x$. Thus the truncated difference is the ordinary finite real
difference, and our added reciprocal-size bound gives
\[
 p_{x+1}+e_s+1/n\leq p_{x+1}+2e_s<p_x.
\]
This implies the required weak additive inequality, including at round zero.
There is no comparison of two different exposures in this argument.

The certificate's witness-thinning, host-thinning, matching and
matching-existence thresholds are at most $N_0\leq n$ for every
$s<Q,x\leq D$. Its matching constant is at most $M\leq C/2$;
division by two is legitimate since $M,C$ are finite. For every ordered
pair $x,y\leq D$, the Four Graphs threshold is at most $N_0\leq n$
and its exponential constant is at least $a$. These are exactly the
table's threshold clauses, including pairs with $x>y$.
Finally its uniform-size clause and static strict reserve give
\[
 n^{-1}\leq\widehat\zeta_s<\zeta_s,\quad
 \widehat\zeta_s+2/n\leq\zeta_s,\quad
 \widehat\zeta_s+1/(\zeta_sn)\leq\zeta_s,
\]
\[
 \omega_{s+1}+2/n\leq e_{s+1},\qquad
 \omega_{s+1}+1/(e_{s+1}n)\leq e_{s+1}.
\]
These verify all six reserve/output premises of the iteration theorem.
The size, degree, host and candidacy premises are the present hypotheses.
The saturating-pair premise is exactly \noderef{SaturatingSourcePair};
nonedges contain no source edges by the reduced-support field of
\noderef{BlowUpInstance}. No additional source degree assumption is needed.

Apply \noderef{VariableSizeInductionIteration}. It returns a colouring
$c:[r]\to[Q]$ and a law $\mu$ on partial compatible embeddings with full
domain, satisfying the invariant at round $D^2=Q-1$ with the entire last
colour class touched. This is a legal completed last round, not an
invariant at round $Q$. The spread conjunct of
\noderef{ColourRoundInvariant} says that this same law's restrictions to
the full domain are $(C/n)$-vertex-spread. Apply
\noderef{EqualSizeTerminalConversion} with domain equal to the full source
set. Its map produces full compatible Copies, agrees pointwise with each
partial map, and preserves the spread assertion under pushforward. These
are the law and bound required in the conclusion; neither conditioning nor
a further loss of mass occurs.
\end{proof}
